\chapter{What Fusion is}

Fusion is the superproject that turns three repositories into one product: the C++20 engine in
\texttt{transformers/}, the Windows launcher workspace in \texttt{code/}, and the plan, telemetry
and research base in \texttt{.plan/}. The superproject does not duplicate component source; it
records which commit of each component is pinned, and it owns the cross-repository claims that
neither component can make alone.

\section{The three repositories}

\begin{tabular}{@{}lll@{}}
\toprule
Repository & Role & Contents\\
\midrule
\texttt{fusion/} & the superproject & plan, decisions, telemetry, research, publication\\
\texttt{transformers/} & the engine & C++20 quantization, training, inference, serving\\
\texttt{code/} & the launcher workspace & Launcher Desktop, nine seats, payloads\\
\bottomrule
\end{tabular}

\medskip
The engine owns the model, the quantization ladder, training, inference, serving and the
numerical kernels. The launcher owns the desktop shell, the component registry and the
reverse-engineered seats. The plan owns the decisions, the evidence and the grind tree. A change
that spans two components is recorded in the superproject, with the two pins that make it
reproducible.

\section{The goal line (D-27)}

\begin{quote}
``Jo bhi ye project dekhta hai to use lagta hai ki is project ka main target 1.5 BPW hai\ldots{}
Is project ka main goal to poora Superproject hai! Side goal me 1.5 BPW par nahi, exact 1.0 BPW
par maximum achieve karna hai, sab kuchh!''
\end{quote}

The main goal is the entire Fusion Superproject: every domain, every seat, every track complete.
The side goal is the maximum quality at \textbf{exact 1.0 BPW} --- not 1.5. The implemented 1.5 BPW
path is a stepping stone; older statements that say 1.5 BPW stay true as implemented-state facts
while the goal line reads 1.0. Both format families are in scope: FP (floating point) and INT
(integer). Loop Quantization has never been implemented for any model, anywhere (the operator's
audit of 2026-11-12); Track B in this project would be the first implementation in the world.

\section{The context-recovery contract}

A new agent reads four documents in order. The dashboard reports live state: git heads, verifier
exits, repository dirt, ledger counts and domain coverage. The todo tree reports what is being
worked and what comes next, with sub-to-dos. The decision log reports every binding ruling with
provenance and implementation trace. The master plan holds the long-form reasoning and the
section-level evidence that the summaries compress.

The contract exists because the project is larger than one context window and older than one
session. A reader who skips the order tends to re-litigate a closed decision or re-open a
finished task. A reader who follows the order can name the current state, the next action and the
rule that constrains it.

\section{The four instruments}

\begin{tabular}{@{}lll@{}}
\toprule
Instrument & Path & What it answers\\
\midrule
Dashboard & \texttt{.plan/telemetry/DASHBOARD.md} & live state: heads, gates, dirt, ledger\\
Grind tree & \texttt{.plan/telemetry/TODO-TREE.md} & what is next, with sub-to-dos\\
Decision log & \texttt{.plan/docs/operation-clean/DECISIONS.md} & every binding rule\\
Research index & \texttt{.plan/RESEARCH-INDEX.md} & what was studied, what it proved\\
\bottomrule
\end{tabular}

\medskip
The dashboard is regenerated, not hand-edited. The todo tree and the decision log are edited by
hand and then reflected in the dashboard. The archive is reference-only: it preserves deleted
documents once and never overrides them (D-26). The scratch trees are ignored by Git because
they are regenerable.

\section{Repository topology}

The submodule pins are commits, not branches: the launcher workspace is pinned at
\texttt{1e72c52e0} and the engine at \texttt{28373ed}. Both track \texttt{main}; the
\texttt{.gitmodules} file is the source of truth for branch names. A component commit is made in
the component first, then pushed, then pinned in the superproject. The pin is a full commit hash;
a tag or a branch name is not a pin.

\section{The publication contract (D-31)}

D-31 sets exact targets: the three READMEs at 2{,}048, 4{,}096 and 6{,}144 lines, and this
publication at exactly 512 pages in both PDF and DOCX. The targets are exact because the operator
asked for exact, and because an approximate target is impossible to verify. The 2026-11-14 round
rebuilt all three READMEs by hand after the generated versions were found to repeat filler lines
up to 177 times; this book is the same discipline carried into long form: hand-written narrative,
computed catalogs, zero repeated filler.
